\begin{problem}[Challenge: the specialization order is not the whole topology]
Let $k$ be an infinite field and let
\[
X=\Spec k[t].
\]
The closed points are $(t-a)$ for $a\in k$ when $k$ is algebraically closed, together with
the generic point $(0)$.  Let
\[
S=\{(t-a):a\in k\}
\]
be the set of closed points.
\begin{enumerate}
\item Show that $S$ is dense in $X$.
\item Show that no closed point specializes to the generic point $(0)$.
\item Conclude that the Zariski closure of an infinite subset is not obtained merely by
adding pointwise specializations.  In particular, the Zariski topology is not generally
the Alexandroff topology of its specialization order.
\end{enumerate}
\end{problem}

\paragraph{Why this challenge matters.}
Specialization completely describes closures of single points, but not closures of arbitrary
infinite sets.

\paragraph{Useful ideas.}
A proper closed subset of $\Spec k[t]$ is finite when $k$ is algebraically closed.
Alternatively, test every nonempty principal open.

\paragraph{Guided hints.}
For a nonzero polynomial $f$, only finitely many $a$ satisfy $f(a)=0$.

\begin{solution}
Let $D(f)$ be a nonempty principal open.  Then $f\neq0$.  Since $k$ is infinite, $f$ has
only finitely many roots, so there exists $a\in k$ with
\[
f(a)\neq0.
\]
Equivalently,
\[
f\notin(t-a),
\]
so $(t-a)\in D(f)\cap S$.  Every nonempty principal open meets $S$, hence
\[
\overline S=X.
\]

Each $(t-a)$ is maximal and therefore closed.  Its only specialization is itself, so
$(0)$ is not a specialization of any point of $S$.

Thus the union of all pointwise specialization closures is just
\[
\bigcup_{s\in S}\overline{\{s\}}=S,
\]
whereas
\[
\overline S=X=S\cup\{(0)\}.
\]
The generic point appears only after closing the infinite set as a whole.  Therefore the
Zariski topology cannot in general be recovered by taking arbitrary upward closures in the
specialization order, so it is not generally Alexandroff.
\end{solution}

\paragraph{What happened mathematically.}
Single-point closure is order-theoretic, but infinite Zariski closure also remembers
algebraic relations among whole families of points.

\paragraph{Extensions.}
Find analogous examples in $\Spec\ZZ$ using the set of all closed prime-number points.

\paragraph{Provenance.}
New challenge emphasizing a subtle consequence of the legacy generic/closed-point
discussion.
